\documentclass[10pt,a4paper]{article}

\usepackage[T1]{fontenc}
\usepackage[utf8]{inputenc}
\usepackage{lmodern}
\usepackage{amsmath,amssymb,amsthm}
\usepackage{booktabs,array,tabularx}
\usepackage{geometry}
\usepackage{xcolor}
\usepackage{enumitem}
\usepackage{microtype}
\usepackage{seqsplit}
\usepackage{hyperref}
\usepackage{fancyhdr}

\geometry{margin=1in}
\hypersetup{
  colorlinks=true,
  linkcolor=blue!55!black,
  citecolor=blue!55!black,
  urlcolor=blue!55!black,
  pdftitle={DSS-MS Unmasked: Public Cross-Join Sanitizer-Key Recovery and Fresh-Message Forgery},
  pdfauthor={Martin Feussner},
  pdfsubject={Cryptanalytic audit report for IACR ePrint 2026/041},
  pdfkeywords={digital signatures, DSS-MS, sanitizable signatures, Chinese remainder theorem, common factor, GCD, key recovery, forgery}
}
\setlist{nosep,leftmargin=*}
\setlength{\parindent}{1em}
\setlength{\parskip}{0.25em}
\setlength{\emergencystretch}{2em}
\Urlmuskip=0mu plus 1mu
\pagestyle{fancy}
\fancyhf{}
\fancyhead[L]{DSS-MS Unmasked}
\fancyhead[R]{18 October 2026}
\fancyfoot[C]{\thepage}
\renewcommand{\headrulewidth}{0.3pt}

\newtheorem{theorem}{Theorem}
\newtheorem{lemma}{Lemma}
\newcommand{\Zq}{\mathbb{Z}_q}
\newcommand{\KeyGen}{\mathsf{KGen}}
\newcommand{\DJoin}{\mathsf{Join}}
\newcommand{\Verify}{\mathsf{Verify}}
\newcommand{\Sanitize}{\mathsf{Sanitizing}}
\newcommand{\hash}[1]{{\ttfamily\small\seqsplit{#1}}}

\title{\bfseries DSS-MS Unmasked:\\
Public Cross-Join Sanitizer-Key Recovery and Fresh-Message Forgery}
\author{Martin Feussner\\
\small Selmer Center, University of Bergen\\
\small \texttt{martin.feussner@uib.no}}
\date{18 October 2026}

\begin{document}
\maketitle
\noindent\fcolorbox{black!55}{black!4}{\parbox{0.94\textwidth}{
This audit used the OpenAI model Daybreak Blue at the ultra reasoning level.
OpenAI Codex generated and selected the paper title and deliberately chose
``Unmasked'' as wordplay on the DSS-MS public sanitizer-key mask revealing the
sanitizer key. OpenAI Codex performed AI-assisted paper and source
reconstruction, cryptanalysis, implementation, experimentation, review
coordination, and technical writing. A separately tasked AI reviewer checked
the sealed primary evidence for internal consistency and returned
\texttt{PASS}. A fresh-context AI agent independently reconstructed the target,
executed eight prospectively sealed full-profile trials, and returned
\texttt{PASS\_INDEPENDENT\_REPRODUCTION}. A separately tasked hostile
falsification review returned \texttt{PASS\_FINAL\_FALSIFICATION\_REVIEW}, and
the final exact-target history recheck returned
\texttt{PASS\_FINAL\_EXACT\_TARGET\_HISTORY\_GATE}. An independent AI
Publication Reviewer reviewed the corrected package and returned
\texttt{PASS\_PUBLICATION\_REVIEW}.
Independent human verification was not performed; it is optional and is not a
publication requirement.}}

\begin{abstract}
DSS-MS, proposed in IACR ePrint 2026/041, publishes a Chinese-remainder
sanitizer-key mask
\(
  S_x=A\,sk_s
\)
as an unreduced integer. The dynamic \(\DJoin\) procedure recomputes the CRT
coefficient after adding sanitizers and publishes
\(S'_x=A'\,sk_s\) under the same global sanitizer public key
\(PK_s=sk_sP\). A passive observer therefore obtains
\[
  \gcd(S_x,S'_x)=sk_s\gcd(A,A').
\]
On the first and only precommitted exact primary instance, containing 3,000
initial sanitizers and one ordinary addition of 20, the residual GCD was one.
One public integer GCD recovered the 251-bit global sanitizer scalar, and the
public equation \(sk_sP=PK_s\) validated the candidate.

Given one valid signature \((R,z,s,T)\) on full message \((m_0,m)\), the
recovered scalar permits the observer to preserve \((R,z)\), retarget the
chameleon-hash component, and obtain a signature accepted on the distinct
full message \((m_0,m')\). Only the sanitizable component changes. The sealed
primary standard-library implementation used actual BLS12-381 G1 arithmetic,
matched the pinned source's SHA-256 transcript conversion, and called a
source-equivalent reconstruction of the ordinary verifier. Its exact run took
34.807209 seconds and 20,480 KiB peak resident memory with one worker. Three
attack-specific negative controls rejected.

An independent implementation ran eight prospectively sealed exact profiles.
Seven recovered exactly after one Join and forged successfully. The retained
eighth trial had residual factor 2 after one Join and at 3,040 sanitizers; its
cumulative public-mask GCD became exact at 3,060 sanitizers, where the forgery
and controls passed. The full eight-trial run took 372.77 seconds and 27,136
KiB peak resident memory with one worker. Thus one-Join recovery is practical
but not guaranteed. The attack recovers the global sanitizer scalar rather
than the original signer scalar, changes only the sanitizable component, and
excludes \(\mathsf{Revoke}\), the separate range-condition correctness defect,
and any compiled-author-binary claim. Repeated-common-factor GCD recovery is
prior art in the PA-CRT lineage; we claim no novelty for the technique. Final
falsification, exact-target history review, and independent Publication Review
passed.
\end{abstract}

\section{Introduction and exact claim}

Zeng, Zhang, Cui, Zhong, and Wang propose a dynamic sanitizable signature with
multiple sanitizers (DSS-MS) for privacy-preserving shared logistics
\cite{dssms}. A signer authenticates a message
\(M=(m_0,m)\), where \(m_0\) is fixed and \(m\) is designated sanitizable.
Every authorized sanitizer receives its own secret \(sk_i\), while the system
uses one global temporary sanitizer scalar \(sk_s\). A public CRT mask lets
each sanitizer derive \(sk_s\) by reduction modulo its own secret. The scheme
adds and revokes sanitizers dynamically.

The addition lifecycle republishes the CRT mask while keeping the same global
scalar. This produces two public integers with a common secret factor. The
exact claim is:

\begin{quote}
For the paper-claimed BLS12-381/SHA-256 profile with 3,000 existing
sanitizers, a passive observer retains the public \(S_x\) values immediately
before and after one ordinary \(\DJoin(20)\). On the precommitted primary
instance, their GCD is exactly \(sk_s\). The observer validates this value
against \(PK_s\), then uses one valid signature to produce an accepted
signature on a fresh full message whose sanitizable component is distinct.
\end{quote}

The result is design-level: both the paper and pinned source use unreduced
integer masks and reuse the scalar through \(\DJoin\). It needs no parser
behavior, malformed point, weak curve, secret channel, attacker-selected
enrollment, rollback, side channel, or implementation memory defect. The
recovered value is not the original signer's Schnorr secret. It is the global
sanitizer scalar, an equivalent signing capability for the scheme's designated
message transformation.

The primary experiment, evidence consistency review, fresh-context independent
reproduction, final falsification review, final exact-target history recheck,
and independent Publication Review passed. The result is classified as a
verified practical design attack under the project charter.

\paragraph{Contributions.}
We identify the common-factor exposure across ordinary public-parameter
epochs; prove the exact key-recovery identity and its public validation rule;
give the fresh-message signing bridge through the scheme's own sanitization
equations; report a precommitted full-profile execution and eight-trial
independent reproduction with controls; state the source and verifier boundary;
distinguish a separate correctness defect; credit the PA-CRT GCD lineage; and
give repair directions.

\section{Target, version, source, and profile}

The target is the current landing-visible body of \emph{Towards
Privacy-Preserving Unmanned Aerial Vehicles Shared Logistics via Dynamic
Sanitizable Signature with Multiple Sanitizers}, IACR ePrint 2026/041
\cite{dssms}. The retained 15-page PDF has SHA-256 digest
\begin{center}
\hash{d35dcd98dc2442b6e263ef58276086832287a3f9d2f9e0131e4b0ad6c7324a20}.
\end{center}
The official landing page records receipt and approval on 11 January 2026. A
version-history endpoint was unavailable during the audit, so this report
does not assign an unsupported revision ordinal.

The paper links the authors' DSS-MS repository \cite{source}. The selected
paper-time source commit and tree are
\begin{center}
\hash{b994a3f46cbe5501f47f53aa2e6942cff714b53a},\\
\hash{5d1b072040170d6df8bca709403d1e01d609536e}.
\end{center}
At the retrieved head
\hash{b48e6359a0c3dce11cc6bf8445577770418d6d27}, the relevant
\texttt{DSS/src/DSS.cpp} remained byte-identical to its paper-time form. The
source fixes 3,000 sanitizers and 256-bit sanitizer values. The paper evaluates
batch additions including 20, and both paper and source accept the addition
count as an input to \(\DJoin\). We therefore freeze profile
\begin{center}
\texttt{DSSMS-BLS12381-SHA256-N3000-ADD20-CRT-INMEM-20260111}.
\end{center}

The source uses BLS12-381 and SHA-256. It provides no fixed portable
known-answer vector or normative wire format. Our endpoint is therefore a
source-equivalent reconstruction of the ordinary equations and transcript
conversion, not execution of a compiled author binary.

\section{DSS-MS equations}

\subsection{Public CRT mask}

Let the sanitizer secrets \(sk_1,\ldots,sk_n\) be pairwise coprime positive
integers. Define
\begin{align}
  N   &= \prod_{i=1}^{n} sk_i, &
  N_i &= N/sk_i, \\
  y_iN_i &\equiv 1 \pmod{sk_i}, &
  A   &= \sum_{i=1}^{n} y_iN_i.                 \label{eq:crt}
\end{align}
The construction samples a global sanitizer scalar \(sk_s\in\Zq^*\), publishes
\begin{equation}
  S_x=A\,sk_s\quad\text{over the integers},
  \qquad PK_s=sk_sP,                              \label{eq:mask}
\end{equation}
and gives sanitizer \(i\) the intended recovery rule
\(sk_s=S_x\bmod sk_i\). Because \(A\equiv1\pmod{sk_i}\), that rule returns
\(sk_s\bmod sk_i\). Literal recovery therefore additionally requires
\(sk_s<sk_i\); Section~\ref{sec:limits} separates this omitted correctness
condition from our attack.

The public nature of \(S_x\) is explicit. \(\KeyGen\) returns it, the paper's
unforgeability experiment gives it to the adversary, the deployment text says
the logistics company publishes it, and the public parameters contain it. The
source stores it as a GMP integer and serializes its complete hexadecimal
representation.

\subsection{Join lifecycle}

Suppose an ordinary \(\DJoin\) adds \(r\) sanitizers. It forms the enlarged
modulus set, recomputes Equation~\eqref{eq:crt} as \(A'\), and publishes
\begin{equation}
  S'_x=A'\,sk_s.                                    \label{eq:joined-mask}
\end{equation}
The named scalar remains \(sk_s\), and \(PK_s\) is not rotated. The pinned
source corroborates the paper: \(\KeyGen\) assigns
\texttt{pp.u\_s = keyPair.sk * u}, while \(\DJoin\) receives the same
sanitizer scalar and assigns \texttt{pp.u\_s = sk\_san * crt.u}. Neither
operation applies a modulus to this multiplication.

\subsection{Signing and verification}

Let \(P\) generate a prime-order group of order \(q\). A signature is
\(\sigma=(R,z,s,T)\). On \(M=(m_0,m)\), the ordinary verifier computes
\begin{align}
  e  &= H_0(m,T), \\
  CH &= T+sP+ePK_s,                                  \label{eq:ch}\\
  c  &= H_1(m_0,R,CH),                               \label{eq:challenge}
\end{align}
and accepts exactly when
\begin{equation}
  zP=R+cPK,                                           \label{eq:verify}
\end{equation}
where \(PK\) is the original signer's public key. An authorized sanitizer
changes \(m\) by preserving \(CH\) and hence the signer's challenge.

\section{Public cross-Join key recovery}

\begin{theorem}[Common-factor recovery]
For public masks \(S_x=A\,sk_s\) and \(S'_x=A'\,sk_s\) computed over the
integers with the same positive scalar \(sk_s\),
\[
  \gcd(S_x,S'_x)=sk_s\gcd(A,A').
\]
If \(\gcd(A,A')=1\), the GCD is exactly \(sk_s\).
\end{theorem}

\begin{proof}
For positive integers, \(\gcd(au,av)=a\gcd(u,v)\). Substitute
\(a=sk_s\), \(u=A\), and \(v=A'\). The coprime case follows immediately.
\end{proof}

The observer need not know any sanitizer secret or either coefficient. It
computes \(g=\gcd(S_x,S'_x)\) directly. The public point supplies an exact
candidate test:
\begin{equation}
  gP\stackrel{?}{=}PK_s.                              \label{eq:public-check}
\end{equation}
If a residual factor \(d=\gcd(A,A')>1\) remains, the one-Join candidate is
\(sk_s d\) and normally fails Equation~\eqref{eq:public-check}. Further public
epochs can reduce the common residual by taking a GCD across all retained
masks. The primary experiment used exactly one Join and had unit residual. The
independent experiment below measured both outcomes: seven of eight sealed
trials had unit residual after one Join, while one retained residual required
two more public epochs. One-Join recovery is therefore not asserted to be
universal.

\section{From the sanitizer scalar to a fresh-message forgery}

\begin{theorem}[Retargeting]
Suppose \((R,z,s,T)\) is accepted on \((m_0,m)\) under
\((PK,PK_s)\), and an attacker knows \(sk_s\) such that
\(PK_s=sk_sP\). For any distinct sanitizable component \(m'\) and any
\(k\in\Zq\), define
\begin{align}
  T' &= CH-kP,                                        \label{eq:tprime}\\
  e' &= H_0(m',T'), \\
  s' &= k-e'sk_s\pmod q.                             \label{eq:sprime}
\end{align}
Then \((R,z,s',T')\) is accepted by the ordinary verifier on the fresh full
message \((m_0,m')\).
\end{theorem}

\begin{proof}
Using \(PK_s=sk_sP\) and Equations~\eqref{eq:tprime}--\eqref{eq:sprime}, the
verifier reconstructs
\begin{align*}
  T'+s'P+e'PK_s
  &= CH-kP+(k-e'sk_s)P+e'sk_sP\\
  &= CH.
\end{align*}
It therefore obtains the same value in Equation~\eqref{eq:challenge}. Since
\((R,z)\) is unchanged, the already valid Equation~\eqref{eq:verify} remains
valid. Because \(m'\ne m\), the full message is distinct.
\end{proof}

This is the scheme's ordinary sanitization transformation performed by a
party that received no sanitizer secret. One valid signature is required. No
sanitizing-oracle query is used. The fresh full message changes only the field
that DSS-MS declares sanitizable; changing \(m_0\) is outside this claim.

\section{Primary full-profile reproduction}

The primary experiment was sealed before execution with decimal seed
\texttt{2026041}. The command fixed \(n=3000\), an addition of 20, 256-bit
sanitizer draws, and 24 explicitly non-full-profile proxy trials. No seed was
discarded or selected after observing the residual GCD.

The dependency-free Python implementation generates source-shaped pairwise
coprime sanitizer values, computes the exact unreduced CRT sums, and performs
the public attack. Its signature layer implements affine arithmetic on actual
BLS12-381 G1 and checks that the generator has prime order. Transcript hashing
matches the pinned source's relevant path: minimal big-endian GMP message
bytes, MIRACL's 49-byte compressed G1 encoding, and a 32-byte SHA-256 digest
loaded into the high end of a zeroed 48-byte integer before reduction modulo
\(q\).

\begin{table}[ht]
\centering
\caption{Sealed primary exact-profile result.}
\label{tab:primary}
\begin{tabularx}{\textwidth}{@{}Xr@{}}
\toprule
Field & Result \\
\midrule
Initial / added / final sanitizers & 3,000 / 20 / 3,020 \\
Old / joined product size & 763,995 / 769,080 bits \\
Recovered public GCD size & 251 bits \\
Residual \(\gcd(A,A')\) & 1 \\
Recovered scalar satisfies \(gP=PK_s\) & yes \\
Honest signature accepted before and after Join & yes \\
Fresh full message distinct & yes \\
Forgery preserves \((R,z)\) and reconstructs the same \(CH\) & yes \\
Forged fresh message accepted & yes \\
Mutation without retargeting accepted & no \\
Wrong-scalar retargeting accepted & no \\
Forgery under wrong \(PK_s\) accepted & no \\
Wall time / peak resident memory / workers & 34.807209 s / 20,480 KiB / 1 \\
\bottomrule
\end{tabularx}
\end{table}

The attack function receives the two public epochs, the public keys, the
public base message and signature, the selected fresh message, and attacker
scalar \(k\). Its interface contains no \(sk_i\) or hidden \(sk_s\) argument.
All tested points lie on the curve and in the prime-order subgroup. The three
rejecting controls distinguish the result from unconditional verifier
acceptance, an irrelevant message mutation, and key-independent retargeting.

The checked-in public package offers two modes. A quick 96-sanitizer mode
exercises the full algebraic, curve, verifier, and control path. The exact mode
replays the deterministic 3,000-sanitizer, \(\DJoin(20)\) primary instance.
Neither is represented as the fresh-context independent gate.

\subsection{Fresh-context independent reproduction}

A fresh-context agent derived a clean-room standard-library implementation
from the frozen paper and source before inspecting the primary reproducer. Its
plan, code, commands, and eight exact 32-byte seeds were prospectively sealed.
Every fixed seed was run once and retained. Each trial began with 3,000
sanitizers and an ordinary \(\DJoin(20)\). A nonunit residual was fixed in
advance to remain a one-Join failure, with at most three further ordinary
public epochs permitted as a separately classified cumulative-GCD test.

\begin{table}[ht]
\centering
\caption{Fresh-context independent full-profile results.}
\label{tab:independent}
\begin{tabular}{@{}crlrl@{}}
\toprule
Trial & One-Join residual & One-Join result & Runtime (s) & Terminal result \\
\midrule
1 & 1 & PASS & 42.548 & PASS one Join \\
2 & 1 & PASS & 42.080 & PASS one Join \\
3 & 1 & PASS & 41.115 & PASS one Join \\
4 & 1 & PASS & 41.226 & PASS one Join \\
5 & 1 & PASS & 41.609 & PASS one Join \\
6 & 2 & retained failure & 70.129 & PASS at 3,060 \\
7 & 1 & PASS & 41.249 & PASS one Join \\
8 & 1 & PASS & 41.929 & PASS one Join \\
\bottomrule
\end{tabular}
\end{table}

All seven one-Join exact recoveries passed Equation~\eqref{eq:public-check},
fresh full-message verification, message-mutation rejection, wrong-key
rejection, honest before/after-Join checks, and curve/subgroup checks. Trial 6
retained residual 2 at 3,020 and 3,040 sanitizers, so no forgery was attempted
with its invalid candidate. The cumulative residual became one at 3,060;
public-key validation, the fresh-message forgery, and all controls then passed.
It remains a one-Join failure and a separately classified multi-epoch success.

The observed one-Join exact-recovery rate is \(7/8\), an empirical result for
fixed seeds rather than an exact probability. The complete eight-trial run,
including the two extra epochs for trial 6, took 6 minutes 12.77 seconds and
27,136 KiB peak resident memory with one worker. The independent manifest has
SHA-256
\begin{center}
\hash{e94ba0b60d291b6ce3136e995f1459dc3fb28147de33908f63cb1c61cb4235e2}.
\end{center}
The independent source, result stream, report, and verdict have respective
SHA-256 digests
\begin{center}
\hash{443b0444f58bf2daf4a0a5951087cf77258a844a49c2b84e2aa115cfacb69aaf},\\
\hash{1a3cd768bcd3106db0f82edc338064faf6118ae574c2b29a6fe8c82a2e087448},\\
\hash{f442e1f684aaf5707bf651e45d5ba2467cd8fecbf3d1b369bff8df6fa960c671},\\
\hash{83d498952d5c88d96202ada5eb52c442ed3a68575323b732dfde76b3267477ca}.
\end{center}
The sealed verdict is \texttt{PASS\_INDEPENDENT\_REPRODUCTION}.

\section{Scope and limitations}\label{sec:limits}

The following qualifications are part of the result.
\begin{itemize}
  \item The primary one-Join attack passively retains two public \(S_x\)
        epochs around one ordinary \(\DJoin(20)\). It does not choose enrolled
        keys or access a secret delivery channel. The independent cumulative
        variant retained further ordinary public epochs only after recording a
        one-Join failure.
  \item The endpoint uses one valid signature. The forged full message changes
        only its sanitizable component, from \((m_0,m)\) to \((m_0,m')\).
  \item The recovered value is the global sanitizer scalar \(sk_s\), not the
        original signer scalar. It nevertheless supplies the intended
        sanitization signing capability to an unauthorized observer.
  \item The quantitative evidence is one precommitted exact primary instance
        plus eight prospectively sealed independent exact trials. Seven of
        eight independent trials recovered after one Join; the sole residual
        trial recovered after two further public epochs. This does not prove an
        exact success probability or a universal epoch bound. The primary
        small proxy trials and random-integer heuristic are not full-profile
        success-rate evidence.
  \item Verification is a source-equivalent reconstruction of the ordinary
        paper and pinned-source equations. The audit did not compile and run
        the author binary and makes no normative codec claim.
  \item The paper and source omit \(sk_s<\min_i sk_i\), so not every honest
        sanitizer necessarily recovers the literal scalar by one reduction.
        That separate correctness defect is excluded from the attack.
  \item The paper and source differ on \(\mathsf{Revoke}\). Revoke is excluded;
        the candidate uses only the paper- and source-consistent
        \(\KeyGen\)-then-\(\DJoin\) lifecycle.
  \item The attack-history result is a bounded search through 18 October 2026,
        not proof of absolute priority. The final exact-target recheck passed;
        private, removed, or unindexed disclosures still cannot be excluded.
\end{itemize}

\section{Prior method and attack history}

This report makes no novelty claim for repeated-common-factor recovery by an
integer GCD. DSS-MS directly cites PA-CRT \cite{pacrt}. Tao et al.'s B-DSPA
work describes public CRT-related products with a reused hidden factor
\cite{bdspa}; Kabil, Aslan, and Azer record that GCD method, analyze its
lifecycle assumptions, and discuss flaws in parts of the preceding PA-CRT
cryptanalysis \cite{kabil}. These works establish relevant technique ancestry.
They do not name DSS-MS, analyze \(S_x=A\,sk_s\), recover the DSS-MS sanitizer
scalar, or reach its ordinary signature verifier.

The bounded exact-target search covered the title, acronym, ePrint identifier,
authors, repository history and issue surfaces, corrections, errata,
sanitizer-key recovery, Join, common factors, GCD, and forgery terms. It
located no earlier public attack against DSS-MS as of 18 October 2026. The
scheme is recent and bibliographic indexing can lag. We therefore claim only
that no prior public exact-target attack was located in the recorded search.
The final recheck returned
\texttt{PASS\_FINAL\_EXACT\_TARGET\_HISTORY\_GATE}. It confirmed that the
official landing and attack-relevant source were unchanged and again found no
exact-target disclosure. The public GCD work against PA-CRT and B-DSPA remains
credited as technique ancestry, and no technique novelty is claimed.

\section{Repair direction}

The public update mechanism must not emit multiple ordinary integers sharing
the same secret scalar as a factor. One design direction is to sample an
independent sanitizer scalar and matching public key for each epoch, bind the
epoch and key into signatures, and give verification explicit version
semantics. Another is to replace the public CRT mask with a conventional
recipient-specific key-establishment mechanism so that public lifecycle
updates do not encode one capability as a reusable common factor.

A revision should also state and enforce the scalar/modulus range needed by
the intended residue recovery, reconcile \(\DJoin\) and \(\mathsf{Revoke}\),
and define a normative transcript and parameter-transition policy. These are
repair directions rather than a validated patch. Any concrete revision needs
fresh correctness, compatibility, and security analysis.

\section{Conclusion}

DSS-MS republishes an unreduced CRT sanitizer-key mask after adding
sanitizers, while retaining the same global sanitizer scalar and public key.
The two public epochs therefore expose that scalar through a standard
common-factor GCD whenever their CRT coefficients are coprime. The one
precommitted exact primary instance had unit residual, recovered the 251-bit
scalar in a 34.807209-second end-to-end run, and used one valid signature to
produce an ordinary-verifier-accepted signature on a fresh full message. The
evidence is limited to the stated Join lifecycle, message transformation, and
source-equivalent verifier boundary. Fresh-context independent reproduction
passed on all eight fixed trials, with seven one-Join successes and one
retained residual that recovered through cumulative public epochs. Final
falsification, exact-target history review, and independent Publication Review
also passed.

\section*{Acknowledgements}
The computations were performed on the Norwegian Research and Education Cloud (NREC), using resources provided by the University of Bergen and the University of Oslo.

\begin{thebibliography}{9}

\bibitem{dssms}
M. Zeng, Q. Zhang, J. Cui, H. Zhong, and F. Wang,
``Towards Privacy-Preserving Unmanned Aerial Vehicles Shared Logistics via
Dynamic Sanitizable Signature with Multiple Sanitizers,''
IACR Cryptology ePrint Archive, Report 2026/041, 2026.
\url{https://eprint.iacr.org/2026/041}

\bibitem{source}
M. Zeng et al.,
``DSS-MS author implementation,'' paper-time commit
\texttt{b994a3f46cbe5501f47f53aa2e6942cff714b53a}.
\url{https://github.com/rookie-papers/DSS-MS}

\bibitem{pacrt}
J. Zhang, J. Cui, H. Zhong, Z. Chen, and L. Liu,
``PA-CRT: Chinese Remainder Theorem Based Conditional Privacy-Preserving
Authentication Scheme in Vehicular Ad-Hoc Networks,''
\emph{IEEE Transactions on Dependable and Secure Computing}, vol.~18, no.~2,
pp.~722--735, 2021.
\url{https://doi.org/10.1109/TDSC.2019.2904274}

\bibitem{bdspa}
Tao et al.,
``B-DSPA: A Blockchain-Based Dynamically Scalable Privacy-Preserving
Authentication Scheme in Vehicular Ad-Hoc Networks,''
\emph{IEEE Internet of Things Journal}, vol.~11, no.~1, p.~1385 ff., 2024.
\url{https://doi.org/10.1109/JIOT.2023.3289057}

\bibitem{kabil}
A. M. Kabil, H. Aslan, and M. Azer,
``Cryptanalysis of Two Conditional Privacy Preserving Authentication Schemes
for Vehicular Ad Hoc Networks,''
\emph{Cryptography}, vol.~8, no.~1, article 4, 2024.
\url{https://doi.org/10.3390/cryptography8010004}

\end{thebibliography}

\end{document}
